%% DRAFT (not yet in v7.tex). New result for Section "Coverage and the Transfer".
%% Uses the paper's macros (\Cold, \Exp, \Hyb, \eps, \R). Verified numerically:
%%   prerun_study/transfer_thm/{sim_full.py,sim_fast.py,dose.py,checks2.py,checks3.py,variant.py}

%% ---------------------------------------------------------------------------
%% (1) Proposed replacement for Proposition "The transfer" (prop:transfer).
%%     The current version defines m from a fixed residual function eps(x) and
%%     treats the rest as noise whose law depends only on M. That misses an
%%     expert that sees information the features do not carry (an LLM reading the
%%     raw text): its choices correlate with the noise. Defining the transfer from
%%     the realized residuals covers both sources, is an exact identity, and is
%%     what the "transfer removed" ablation of Section 4 actually removes.
%% ---------------------------------------------------------------------------

\begin{proposition}[The transfer]
\label{prop:transfer}
For any reference $\theta^\star\in\R^d$, a warm start built from
$(x_s,r_s)_{s\le\tau}$ satisfies, exactly,
\[
\widehat\theta_1=\theta^\star-\lambda V_1^{-1}\theta^\star+m,
\qquad
m:=V_1^{-1}\sum_{s\le\tau}\bigl(r_s-\langle\theta^\star,x_s\rangle\bigr)\,x_s .
\]
Given $M$, the \emph{transfer} $m$ is the only term that depends on which
rounds were logged and what they returned. Under the linear model, with noise
independent of the routed actions, $m=V_1^{-1}\xi$ and its law is fixed by $M$
(Proposition~\ref{prop:response-sum}). It departs from that law in two ways:
through a residual $\varepsilon(x)=\E[r\mid x]-\langle\theta^\star,x\rangle$
(misspecification), and through an expert whose choices depend on information
about the reward that the features do not carry, as an LLM reading the raw
text does; the second makes
$\E[r_s-\langle\theta^\star,x_s\rangle\mid x_s]\ne0$ on the routed rounds even
when $\varepsilon\equiv0$.
\end{proposition}

%% ---------------------------------------------------------------------------
%% (2) New result.
%% ---------------------------------------------------------------------------

\paragraph{Equal coverage, unequal value.}
Proposition~\ref{prop:response-sum} says that under linear rewards two experts
of equal coverage are interchangeable. The next result shows how far this
fails once the expert sees more than the learner. It uses the canonical
instance of Proposition~\ref{prop:three-way} with one change: the good action
is right only on a fraction $p$ of the rounds, and which rounds is visible to
the expert but not to the learner.

\emph{The instance.} Fix $k\ge1$, $0\le q<p$ with $p+q\le1$, $L>1$,
$\lambda>0$, a horizon $T>4k$, and $d=T+2$. At round $t$ the learner is
offered $D=e_1$, $R=e_2$ and a fresh arm $b_t=Le_{t+2}$. Each round has a type
$y_t\in\{D,R,\varnothing\}$ that the learner does not see: the action named by
$y_t$ has reward $1$ and every other action reward $0$ (so $b_t$ is never
rewarded). The first $\tau:=4k$ rounds contain $k$ rounds of type $D$, $k$ of
type $R$ and $2k$ of type $\varnothing$, in any order; later types are
i.i.d.\ with $\Prb(D)=p$, $\Prb(R)=q$. The learner's expected reward is linear
in its features, $\mu(x)=\langle\theta^\star,x\rangle$ with
$\theta^\star=pe_1+qe_2$, and regret is measured against it: the gaps are
$0$ for $D$, $p-q$ for $R$ and $p$ for every $b_t$. Two experts see the type.
$\Exp_A$ plays $D$ on rounds of types $D$ and $R$ and plays $R$ on rounds of
type $\varnothing$; $\Exp_B$ plays $R$ on rounds of types $D$ and $R$ and plays
$D$ on rounds of type $\varnothing$. LinUCB uses regularization $\lambda$ and a
constant radius $\beta>0$; we write $c:=\beta L/\sqrt\lambda$ for the index of
an untried fresh arm, and break ties in favour of the fresh arm.

\begin{theorem}[Equal coverage, unequal value]
\label{thm:transfer-separation}
On this instance the following hold.
\begin{enumerate}[itemsep=1pt,topsep=1pt,leftmargin=*]
\item[(i)] \emph{(Same coverage, same cost, same accuracy.)}
$M_\tau^{A}=M_\tau^{B}=2k(e_1e_1^\top+e_2e_2^\top)$, both acquisition regrets
equal $2k(p-q)$, and each expert is right on exactly $k$ of its $\tau$ routed
rounds. The two hybrids therefore hand off with the same design matrix, hence
the same widths, certified directions and prices, and their estimates differ
only through the transfer:
$\widehat\theta^{A}_{\tau+1}=\frac{k}{\lambda+2k}\,e_1$ and
$\widehat\theta^{B}_{\tau+1}=\frac{k}{\lambda+2k}\,e_2$.
\item[(ii)] \emph{(Cold, and coverage alone.)} Cold LinUCB plays the fresh arm
at every round, so $R_T(\Cold)=pT$. The same holds after round $\tau$ for
either hybrid if the routed rounds contribute their covariance but not their
rewards.
\item[(iii)] \emph{(The expert whose right answers are rare.)}
$\Hyb^{B}_\tau$ never plays $D$ after round $\tau$, so
$R_T(\Hyb^{B}_\tau)\ge(p-q)(T-2k)$.
\item[(iv)] \emph{(The expert whose right answers are frequent.)} Suppose
$c\le\min\{p,\tfrac12\}/2$ and $h:=k(1-2c)-c\lambda>0$. With probability at
least
\[
1-e^{-8(p-c)h}\ \ \ge\ \ 1-e^{-2p(k-2c\lambda)},
\]
$\Hyb^{A}_\tau$ plays $D$ at every round after $\tau$, and then
$R_T(\Hyb^{A}_\tau)=2k(p-q)$.
\end{enumerate}
Consequently, with that probability,
\[
\frac{R_T(\Cold)}{R_T(\Hyb^{A}_\tau)}=\frac{pT}{2k(p-q)},
\qquad\text{while}\qquad
\frac{R_T(\Cold)}{R_T(\Hyb^{B}_\tau)}\le\frac{p}{p-q}\cdot\frac{T}{T-2k}.
\]
\end{theorem}

Two experts that buy exactly the same coverage, at exactly the same price,
with exactly as many right answers, give warm starts whose multipliers over
cold LinUCB differ by a factor of order $T/k$. Nothing that depends on
coverage alone---the design matrix, the widths, the prices, the price bound of
Theorem~\ref{thm:general-price}---can tell them apart, and neither can the
experts' accuracy. What separates them is where their right answers fall:
$\Exp_A$'s are on an action that is right on a fraction $p$ of all later
rounds, $\Exp_B$'s on one that is right on a fraction $q$. The number of right
answers needed does not grow with the horizon: the failure probability in (iv)
decays exponentially in $k$ alone.

\begin{remark}[Where the theorem sits]
\label{rem:transfer-regime}
The instance violates the hypotheses of both earlier results, and each
violation is needed. The experts see the type, so their logged rewards are not
independent of their actions, which Proposition~\ref{prop:response-sum}
assumes. And $c\le p/2$ forces $\beta<\sqrt\lambda\,\|\theta^\star\|$, below
any valid confidence radius, which Theorem~\ref{thm:general-price} assumes.
If instead $\beta(L-1)/\sqrt\lambda\ge1$, the regime of
Proposition~\ref{prop:one-pass}, then $\langle\widehat\theta_t,D\rangle<1$
gives $\mathrm{UCB}_t(D)<1+\beta/\sqrt\lambda\le c$ for every warm start, so no
hybrid plays $D$ after $\tau$ and neither expert helps. With a valid radius,
coverage governs; with the annealed exploration used in practice
(Section~\ref{sec:exp-setup}), the transfer decides.
\end{remark}

\begin{remark}[Settling, and a more accurate but useless expert]
\label{rem:transfer-extensions}
(a) The theorem also shows the all-or-nothing behaviour seen on AmazonCat
(Appendix~\ref{app:dominant}): once $\Hyb^A_\tau$ plays a fresh arm, the
statistics of $D$ stop changing and every later fresh arm has index $c$, so it
never plays $D$ again. It either settles on $D$ for good or never does.
(b) $\Exp_B$ can be made strictly more accurate without changing anything
after $\tau$: add $j$ window rounds on which the round's fresh arm is right and
$j$ on which no action is; let $\Exp_B$ play the fresh arm on the former and
$D$ on the latter, and $\Exp_A$ the reverse. Both still play $D$ equally often
and differ only on fresh directions that never recur, $\Exp_B$ is right on
$k+j$ routed rounds against $\Exp_A$'s $k$. Both acquisition regrets become
$2k(p-q)+jp$, each hybrid has played $D$ $2k+j$ times at handoff, and
(iii)--(iv) go through with $2k+j$ in place of $2k$ in the count of $D$, so
that (iv) needs $h:=k-c(\lambda+2k+j)>0$.
(c) The proof uses the radius only through $c$, so the theorem holds for any
non-increasing schedule $\beta_t$ with $c:=\sup_{t>\tau}\beta_tL/\sqrt\lambda$.
\end{remark}

%% ---------------------------------------------------------------------------
%% (3) Proof (for the appendix).
%% ---------------------------------------------------------------------------

\begin{proof}[Proof of Theorem~\ref{thm:transfer-separation}]
Every offered arm is a multiple of a coordinate vector, so for every policy
$V_t$ is diagonal and each coordinate of $\widehat\theta_t$ is computed
separately. If before round $t$ the action $D$ has been played $n_D$ times with
total reward $s_D$, and $R$ has been played $n_R$ times with total reward
$s_R$, then
\[
\mathrm{UCB}_t(D)=\frac{s_D}{\lambda+n_D}+\frac{\beta}{\sqrt{\lambda+n_D}},\qquad
\mathrm{UCB}_t(R)=\frac{s_R}{\lambda+n_R}+\frac{\beta}{\sqrt{\lambda+n_R}},\qquad
\mathrm{UCB}_t(b_t)=c,
\]
the last because coordinate $t+2$ has not been touched before round $t$.
Since $L>1$, $c=\beta L/\sqrt\lambda>\beta/\sqrt{\lambda+n}$ for every
$n\ge0$.

\emph{(i)} On the window, $\Exp_A$ plays $D$ on the $k$ rounds of type $D$
(right) and the $k$ of type $R$ (wrong), and $R$ on the $2k$ of type
$\varnothing$ (wrong); $\Exp_B$ plays $R$ on the $k$ rounds of type $D$ (wrong)
and the $k$ of type $R$ (right), and $D$ on the $2k$ of type $\varnothing$
(wrong). Each plays $D$ and $R$ exactly $2k$ times, so
$M^A_\tau=M^B_\tau=2k(e_1e_1^\top+e_2e_2^\top)$, each pays $2k(p-q)$ for its
plays of $R$, and each is right $k$ times. The response sums are $ke_1$ and
$ke_2$, which gives the two estimates.

\emph{(ii)} Cold LinUCB has $n_D=n_R=0$, so $\mathrm{UCB}_t(D)=\mathrm{UCB}_t(R)
=\beta/\sqrt\lambda<c$ and it plays $b_t$; by induction it does so at every
round, paying $p$ each time. If the routed rounds contribute no rewards, then
after $\tau$ we have $s_D=s_R=0$, both indices equal $\beta/\sqrt{\lambda+2k}<c$,
and the same induction applies.

\emph{(iii)} After the window $\Hyb^B_\tau$ has $n_D=2k$ and $s_D=0$, because
$\Exp_B$ played $D$ only on rounds of type $\varnothing$. Hence
$\mathrm{UCB}_t(D)=\beta/\sqrt{\lambda+2k}<c$, and as long as $D$ is not
played these statistics do not change; so $D$ is never played. Every round
after $\tau$ then costs at least $p-q$, and the window cost $2k(p-q)$, so
$R_T(\Hyb^B_\tau)\ge2k(p-q)+(T-4k)(p-q)=(p-q)(T-2k)$.

\emph{(iv)} Let $X_i:=\ind\{y_{\tau+i}=D\}$, which are i.i.d.\ Bernoulli$(p)$,
$S_n:=\sum_{i\le n}X_i$, and
$G:=\{k+S_n>c(\lambda+2k+n)\ \text{for all }n\ge0\}$. On $G$, $\Hyb^A_\tau$
plays $D$ at every round after $\tau$: if it has done so at rounds
$\tau+1,\dots,\tau+n$, then $n_D=2k+n$ and $s_D=k+S_n$, so
$\mathrm{UCB}(D)>s_D/(\lambda+n_D)>c$; and $n_R=2k$, $s_R=0$, so
$\mathrm{UCB}(R)=\beta/\sqrt{\lambda+2k}<c$. It therefore plays $D$ at round
$\tau+n+1$. Every round after $\tau$ then has gap $0$, and
$R_T(\Hyb^A_\tau)=2k(p-q)$.

For the probability, $n=0$ is in $G$ because $h>0$. For $n\ge1$, put
$Y_i:=c-X_i\in[c-1,c]$, so that $\E Y_i=c-p<0$ and $G$ fails exactly when
$\sum_{i\le n}Y_i\ge k-c(\lambda+2k)=h$ for some $n\ge1$. By Hoeffding's lemma,
$\E e^{\theta Y_i}\le\exp(\theta(c-p)+\theta^2/8)$, which equals $1$ at
$\theta=8(p-c)>0$. So $\exp(\theta\sum_{i\le n}Y_i)$ is a nonnegative
supermartingale starting at $1$, and Ville's maximal inequality gives
$\Prb(G^c)\le e^{-\theta h}=e^{-8(p-c)h}$. Finally $c\le p/2$ gives
$8(p-c)\ge4p$ and $c\le1/4$ gives $h\ge(k-2c\lambda)/2$, so
$8(p-c)h\ge2p(k-2c\lambda)$. The ratios follow from (ii)--(iv).
\end{proof}
